\section{Statement, relation, and state transition}
\label{sec:relation}

This section fixes the object proved by \Profile{}.  The distinction between
the arithmetic relation and the on-chain state preconditions is important:
the proof claim is satisfiability of the former, while the
accepting instruction checks the latter directly before changing any account.

\subsection{Fields, digests, and domain-separated hashes}

Let
\[
  p=2^{31}-1,\qquad \mathbb F=\mathbb F_p,
\]
and fix the tower
\[
  \mathbb F_1=\mathbb F[i]/(i^2+1),\qquad
  \mathbb K=\mathbb F_1[u]/(u^2-(2+i)).
\]
Thus \(|\mathbb K|=p^4\); the implementation denotes this fixed
quadratic-over-quadratic representation by \textsf{QM31}.  A payment digest is an element of
\(\mathbb F^8\).  All field elements have canonical little-endian encodings;
an encoding of an integer greater than or equal to \(p\) is rejected.

The arithmetic relation uses a fixed width-16 Poseidon2 permutation over
\(\mathbb F\)~\cite{GrassiKhovratovichSchofnegger2023}.  We write
\(H_{\mathsf{own}}\), \(H_{\mathsf{nul}}\), and \(H_{\mathsf{note}}\) for
the corresponding domain-separated sponge calls, and
\(H_{\mathsf{node}}\) for the width-16 two-to-one compression used by the
atomic payment tree.  The domains are part of the relation, rather than
metadata supplied by a prover.  In particular, an output note and a future
input note use the same \(H_{\mathsf{note}}\) domain.  For digests
\(L,R\in\mathbb F^8\), the atomic node compressor places \(L\|R\) in the
permutation state, adds the fixed atomic-v3 tweak in the final lane, applies
the permutation, and returns the first eight lanes.

For a bit \(b\), define
\[
  \operatorname{Parent}_b(X,S)=
  \begin{cases}
     H_{\mathsf{node}}(X,S),&b=0,\\
     H_{\mathsf{node}}(S,X),&b=1.
  \end{cases}
\]
For \(i=\sum_{j=0}^{19}i_j2^j\) and
\(\boldsymbol\sigma=(\sigma_0,\ldots,\sigma_{19})\), set
\[
  \operatorname{Root}(L;i,\boldsymbol\sigma)
   =X_{20},\quad X_0=L,\quad
   X_{j+1}=\operatorname{Parent}_{i_j}(X_j,\sigma_j).
\]

\subsection{Public statement and witness}

\begin{definition}[\Profile{} public statement]
\label{def:spend-statement}
A public statement is
\[
 x=(P,s,A,\nu,C_{\mathrm{out}},A',a,f,\Delta),
\]
where \(P\in\{0,1\}^{256}\) is the pool account address,
\(s\in\{0,\ldots,2^{64}-1\}\) is its sequence number,
\(A,\nu,C_{\mathrm{out}},A'\in\mathbb F^8\),
\(a\in\mathbb F\) is an asset identifier,
\(f\in\{0,\ldots,2^{30}-1\}\) is the public fee, and
\(\Delta\in\{0,1\}^{256}\) is the pool's deployment domain.
The canonical 216-byte encoding is
\[
 \mathtt{version}\|\mathtt{depth}\|0^6\|P\|s_{\rm le}\|A\|\nu\|
 C_{\mathrm{out}}\|A'\|a_{\rm le}\|f_{\rm le}\|\Delta,
\]
with version \(4\), depth \(20\), 32 bytes per digest, and four bytes for
each of \(a\) and \(f\).  SHA-256 with the fixed statement domain hashes
this encoding into the Fiat--Shamir transcript.
\end{definition}

The deployment domain binds each statement to one program deployment on one
named cluster.  At pool initialization the program computes
\[
 \Delta=\operatorname{SHA-256}(
   \mathtt{aspis\mbox{-}spend\mbox{-}deployment\mbox{-}domain\mbox{-}v1}
   \,\|\,\mathrm{pid}\,\|\,\tau),
\]
where \(\mathrm{pid}\) is the runtime program id observed by the executing
instruction and \(\tau\) is a nonempty domain tag of at most 64 bytes
(for example \texttt{mainnet-beta} or \texttt{devnet}) supplied on the
initialization wire, and stores \(\Delta\) in the pool account.  Every
state-affecting verification requires the statement's \(\Delta\) to equal
the pool's stored value before any other precondition involving proof bytes
is evaluated.  Because the statement digest is already absorbed into the
Fiat--Shamir transcript, \(\Delta\) is transcript-bound without a separate
absorb step: a proof ground for one deployment domain cannot satisfy the
transcript schedule of any statement carrying another.

\begin{definition}[\Profile{} witness]
\label{def:spend-witness}
A witness is
\[
 w=(k_\nu,r_{\mathrm{in}},r_{\mathrm{out}},pk_{\mathrm{out}},
    a_{\mathrm{in}},v,v',i,
    \sigma_0,\ldots,\sigma_{19},
    \boldsymbol\ell_{\mathrm{in}},
    \boldsymbol\ell_{\mathrm{out}}),
\]
where each of \(k_\nu\), \(r_{\mathrm{in}}\),
\(r_{\mathrm{out}}\), \(pk_{\mathrm{out}}\), and
\(\sigma_0,\ldots,\sigma_{19}\) lies in \(\mathbb F^8\),
\(a_{\mathrm{in}}\in\mathbb F\),
\(v,v'\in\{0,\ldots,2^{30}-1\}\),
\(i\in\{0,\ldots,2^{20}-1\}\), and each range vector belongs to
\(\{0,\ldots,2^{10}-1\}^3\).
\end{definition}

\begin{definition}[Atomic same-path relation]
\label{def:spend-relation}
The relation \(\RProfile(x,w)=1\) holds exactly when all of the following
equations hold over the indicated integer or field domain:
\begin{align}
 pk_{\mathrm{in}}
   &=H_{\mathsf{own}}(k_\nu),                                                   \label{eq:owner-key}\\
 L_{\mathrm{in}}
   &=H_{\mathsf{note}}(pk_{\mathrm{in}},v,a_{\mathrm{in}},r_{\mathrm{in}}),    \label{eq:input-note}\\
 \nu
   &=H_{\mathsf{nul}}(k_\nu,r_{\mathrm{in}}),                                  \label{eq:nullifier}\\
 C_{\mathrm{out}}
   &=H_{\mathsf{note}}(pk_{\mathrm{out}},v',a,r_{\mathrm{out}}),               \label{eq:output-note}\\
 A
   &=\operatorname{Root}(L_{\mathrm{in}};i,\boldsymbol\sigma),                  \label{eq:input-root}\\
 A'
   &=\operatorname{Root}(C_{\mathrm{out}};i,\boldsymbol\sigma),                \label{eq:output-root}\\
 a_{\mathrm{in}}&=a,                                                            \label{eq:asset-equality}\\
 v' + f&=v,                                                                      \label{eq:balance}\tag{integer}\\
 v&=\ell_{\mathrm{in},0}+2^{10}\ell_{\mathrm{in},1}
       +2^{20}\ell_{\mathrm{in},2},                                             \label{eq:range-in}\\
 v'&=\ell_{\mathrm{out},0}+2^{10}\ell_{\mathrm{out},1}
       +2^{20}\ell_{\mathrm{out},2}.                                            \label{eq:range-out}
\end{align}
Each of the six limbs in \eqref{eq:range-in}--\eqref{eq:range-out} is
proved to be a member of the exact 10-bit lookup table.  Equation
\eqref{eq:balance} is an integer equality after the 30-bit range checks; it
is not an equality modulo \(p\).
\end{definition}

The same \(i\) and the same twenty siblings occur in
\eqref{eq:input-root} and \eqref{eq:output-root}.  Thus the relation replaces
one private leaf along one private path.  It does not prove a public append,
a multi-input or multi-output transfer, or a statement about an otherwise
general shielded pool.  The fields \(P\) and \(s\) are transcript-bound public
context.  Their correspondence to live state is checked by the instruction,
not inferred from the Merkle equations; in particular, \(s\) is a pool
revision and is not the private path index \(i\).

\subsection{Trace realization}

The arithmetic trace has \(2^{10}=1024\) rows.  Its 49 Poseidon blocks comprise
one owner-key block, three input-note sponge blocks, twenty input-path
compressions, twenty output-path compressions, two nullifier sponge blocks,
and three output-note sponge blocks.  The two
path computations share the physical direction-bit and sibling sources at
every level.  Tagged copy constraints enforce this sharing, while the
Poseidon2, lookup, balance, public-binding, and terminal constraints enforce
Equations~\eqref{eq:owner-key}--\eqref{eq:range-out}.  The generated atomic-v3
registry contains 183 copy terms and fixes the trace layout independently of
the witness.

\begin{proposition}[Normalized trace data implies the spend relation]
\label{prop:relation-realization}
Assume that the deployed Poseidon2 constants, domain labels, absorption rules,
and output truncation equal the functions in
Equations~\eqref{eq:owner-key}--\eqref{eq:output-note}.  Let normalized trace
data for a statement \(x\) contain:
\begin{enumerate}[label=(\roman*)]
  \item zero arithmetic residuals for both range checks, asset equality, and
  balance;
  \item the required typed owner, note, nullifier, and Merkle-node input for
  every permutation, its eleven checked two-round transitions, and its output,
  including the child-selection and node-absorption equations at each Merkle
  level;
  \item twenty input-path and twenty output-path levels using the same path
  bits and siblings; and
  \item openings equal to the six public spend fields in \(x\).
\end{enumerate}
Then those data determine integers \(v,v'\) and a witness \(w\) such that
\(\RProfile(x,w)=1\).  Conversely, every \((x,w)\) satisfying
\(\RProfile(x,w)=1\) can be encoded as normalized trace data with all of those
residuals zero.
\end{proposition}

\begin{proof}
For the first direction, zero arithmetic residuals give the two 30-bit range
checks, balance, and asset equality.  The eleven two-round transitions expand
to the specified twenty-two Poseidon2 rounds.  The owner, note, nullifier, and
output-note blocks therefore give the four domain-separated hash equations.
At each Merkle level, the two selection residuals give the ordered pair
\((X_j,\sigma_j)\) or \((\sigma_j,X_j)\) according to the shared Boolean path
bit.  Chaining the twenty levels and using the two terminal root equalities gives
\eqref{eq:input-root} and \eqref{eq:output-root}.  The Boolean-bit constraints
and limb/value reconstruction equations give the 30-bit ranges, after which the
balance constraint is the integer equality \eqref{eq:balance}.  The remaining
public terminals enforce \eqref{eq:nullifier}, \eqref{eq:output-note}, and
\eqref{eq:asset-equality}.  These values form the required witness.

For the converse, given a witness satisfying the displayed equations, fill each
hash block with the deterministic Poseidon2 execution, place the common path
bits and siblings in their registry cells, and use the stated 10-bit limbs.
Every arithmetic, hash, path, copy, and public-field residual in the normalized
data then vanishes.  This
argument concerns the normalized data only.  It does not show that arbitrary
accepted Rust execution yields those data.  That separate step must justify
the byte decoding, absorption-row substitutions, copy and LogUp sharing,
polynomial-commitment and FRI extraction, transcript argument, and public and
live-account bindings.  Section~\ref{sec:soundness} records that step as an
explicit premise, and the Lean theorem
\texttt{AspisV5AcceptedSpendRelation.extracted\_trace\_implies\_spend\_relation}
checks the deterministic implication above.
\end{proof}

\begin{lemma}[No repeated public nullifier]
\label{lem:no-double-spend}
The transition checks the program-derived account at seeds
\(\mathtt{aspis\mbox{-}nullifier\mbox{-}v1}\Vert\nu\), records that account as
spent atomically on acceptance, and rejects any later transition carrying the
same public nullifier \(\nu\).  Subject to Solana's account-locking and atomic
execution semantics, at most one successful state transition can therefore
carry a given public nullifier.

This program fact is narrower than a cryptographic theorem saying that one
semantic note can never be opened under two different nullifiers.  For a fixed
note leaf
\[
 L=H_{\mathsf{note}}(H_{\mathsf{own}}(k_\nu),v,a,r_{\mathrm{in}}),
\]
the honest opening determines
\(\nu=H_{\mathsf{nul}}(k_\nu,r_{\mathrm{in}})\).  An accepted second spend of
that same leaf with a different nullifier would require either failure of the
proof's extraction and relation binding, a different opening of the combined
owner-and-note commitment, or a Merkle binding failure.  It is not a
target-second-preimage event for \(H_{\mathsf{nul}}\), because a second
preimage has the same hash output.  The target-second-preimage reduction in
Corollary~\ref{cor:theft-resistance} instead concerns a wrong extracted secret
for a fixed public nullifier.  The Lean theorem
\texttt{different\_opening\_same\_leaf\_reduction} checks the logical
fixed-leaf note-opening reduction.  It does not prove that an alternative leaf
or path reaches the same root, and this paper assigns no numerical bound to
either commitment event.
\end{lemma}

\subsection{Live-state preconditions and atomic effect}

For an accepting state-changing instruction, the verifier additionally
requires the following public pre-state: (i) the supplied pool account has
the expected owner, type, 80-byte length, key \(P\), sequence
\(s<2^{64}-1\), anchor
\(A\), and stored deployment domain equal to the statement's \(\Delta\),
with a distinct error code for a domain mismatch checked before the anchor
comparison; (ii) the proof account has the expected owner and a finalized
  all-zero authority sentinel; and (iii) the nullifier marker is the expected
  program-derived address with seeds
\(\mathtt{aspis\mbox{-}nullifier\mbox{-}v1}\) and \(\nu\), is unspent, and
has one of the accepted pre-state shapes.  Its 72-byte post-state stores and
checks both \(P\) and \(\nu\).  The instruction
constructs \(x\) from these accounts and its public arguments, verifies the
entire proof, reborrows and rechecks writable state, and only then commits
\[
   (s,A)\longmapsto(s+1,A')
\]
and the marker for \((P,\nu)\).  A failure before the final commit leaves both
objects unchanged.  These account checks prevent a valid proof for one pool,
revision, anchor, deployment domain, or nullifier from authorizing a
different transition.

\begin{lemma}[Deployment-domain separation]
\label{lem:deployment-domain}
Let a pool store deployment domain \(\Delta'\) and let a proof carry a
statement with deployment domain \(\Delta=\mathrm{sha256}(
\mathtt{aspis\mbox{-}spend\mbox{-}deployment\mbox{-}domain\mbox{-}v1}
\Vert\mathit{program\_id}\Vert\mathit{domain\_tag})\)
(Definition~\ref{def:spend-statement}).  If \(\Delta\neq\Delta'\) the
transition is rejected before the proof is verified and before any account is
mutated, except with negligible probability.  \(\Delta\) is part of the
statement-digest preimage, absorbed under the \textsf{Statement} label, so by
Theorem~\ref{thm:argument-of-knowledge} an accepting proof is sound only for
the \(\Delta\) in its own statement; independently the program compares the
statement's \(\Delta\) against the pool's stored \(\Delta'\) fail-closed and
returns a dedicated mismatch error before verification and before any mutable
borrow.  The guarantee is separation between pools with distinct stored
domains, not cross-cluster unforgeability: because \(\mathit{domain\_tag}\) is
operator-asserted and not bound to cluster genesis, a keyholder can reproduce
an identical \(\Delta\) elsewhere with the same \(\mathit{program\_id}\) and
tag.
\end{lemma}
